% TOP_PROBLEM: {"id": 3373, "title": "Are there only finite Fermat Primes?", "category_id": 1, "category": {"id": 1, "name": "number_theory", "display_name": "Number Theory", "description": "Properties of integers, prime numbers, Diophantine equations.", "slug": "number-theory", "order_index": 1, "created_at": "2026-07-31T15:26:25.670Z"}, "status": "open"}
% FIRST_POSTED: 2026-09-25
% ENTRY_KIND: statement_only
% !TeX program = lualatex
% Definition and sources only; this file is not an LLM solution attempt.
\documentclass[11pt]{article}
\usepackage[margin=25mm]{geometry}
\usepackage{fontspec}
\setmainfont{Latin Modern Roman}
\usepackage{amsmath,amssymb,mathtools}
\usepackage{unicode-math}
\setmathfont{Latin Modern Math}
\usepackage{xurl,hyperref}
\hypersetup{colorlinks=true,urlcolor=blue}
\setcounter{secnumdepth}{0}
\setlength{\parindent}{0pt}
\setlength{\parskip}{5pt}
\setlength{\emergencystretch}{3em}
\begin{document}
\section*{\texorpdfstring{Are there only finite Fermat Primes?}{Are there only finite Fermat Primes?}}\label{3373}
\subsection{Definitions and mathematical statement}
For each integer $n\ge0$, the Fermat number is
\[
 F_n=2^{2^n}+1.
\]
A Fermat prime is a member of this sequence with no positive divisors other than one and itself. Determine whether the index set
\[
 \{n\in{\mathbb{Z}}_{\ge0}:F_n\text{ is prime}\}
\]
is infinite or finite. This is a binary determination problem, not a requirement to prove infinitude in advance. Proving compositeness for any finite set of further terms does not establish eventual compositeness, and hence does not by itself establish finiteness of the prime terms.
\par\smallskip\textit{Formulation and concept references:} \href{https://webusers.imj-prg.fr/~michel.waldschmidt/articles/pdf/NTchallengesVI.pdf}{[S1]}.

\subsection{Short English statement}
Are there infinitely many prime Fermat numbers, or are only finitely many terms of this specific sequence prime?
\subsection{Sources}
\begin{itemize}
\item[S1]Michel Waldschmidt, "Number theory: Challenges of the twenty-first century," Delhi University survey lecture, 12 October 2012.
\url{https://webusers.imj-prg.fr/~michel.waldschmidt/articles/pdf/NTchallengesVI.pdf}.
\textit{Formulation source.}
\end{itemize}
\end{document}
